% CONTEXT REFERENCE ARCHIVE -- NOT AN ACTIVE BIBLIOGRAPHY INPUT.
% These individually reviewed entries were moved from the maintained bibliography
% on 2026-08-31 because the current paper makes no specific claim citing them.
% Entry text and stable keys are preserved verbatim; no reference record is lost.
% Neither maintained driver inputs this file. Restore an entry to the active
% bibliography only with a supporting citation, then run bibliography QA.
% The separate conditional auditable-budget records remain in their original file.

% Existing selective/reject-option sources support the active abstention contrast.
\KBbibitem{kalai2021abstain}{Kalai and Kanade}{2021} A. T. Kalai, V. Kanade. Towards optimally abstaining from
prediction with OOD test examples. \emph{NeurIPS}, 2021. arXiv:2105.14119.

% The active dataset table cites ImageNet-C/R sources, not ImageNet construction.
\KBbibitem{deng2009imagenet}{Deng et al.}{2009} J. Deng, W. Dong, R. Socher, L.-J. Li, K. Li, L. Fei-Fei.
ImageNet: A large-scale hierarchical image database. \emph{CVPR}, 2009.

% No DomainBed implementation or comparison is described in the maintained paper.
\KBbibitem{gulrajani2021domainbed}{Gulrajani and Lopez-Paz}{2021} I. Gulrajani, D. Lopez-Paz. In search of lost
domain generalization. \emph{ICLR}, 2021.

% No active RobustBench claim or identified dependency in the checked canonical
% metadata. Revisit this citation if released-checkpoint attribution evidence emerges.
\KBbibitem{croce2021robustbench}{Croce et al.}{2021} F. Croce, M. Andriushchenko, V. Sehwag,
E. Debenedetti, N. Flammarion, M. Chiang, P. Mittal, M. Hein.
RobustBench: A standardized adversarial robustness benchmark.
\emph{NeurIPS Datasets and Benchmarks}, 2021.

% General context is covered by the specific adapter and monitoring sources.
\KBbibitem{liang2025ttasurvey}{Liang et al.}{2025} J. Liang, R. He, T. Tan. A comprehensive survey on
test-time adaptation under distribution shifts. \emph{IJCV} 133(1):31--64, 2025.
arXiv:2303.15361.

% The paper does not invoke a PAC covariate-shift prediction-set theorem.
\KBbibitem{park2022pac}{Park et al.}{2022} S. Park, E. Dobriban, I. Lee, O. Bastani. PAC prediction sets
under covariate shift. \emph{ICLR}, 2022. arXiv:2106.09848.

% No PeTTA recurring-shift comparison is made in the maintained paper.
\KBbibitem{hoang2024petta}{Hoang et al.}{2024} T. Hoang, D. M. Vo, M. N. Do. Persistent test-time adaptation
in recurring testing scenarios. \emph{NeurIPS}, 2024. arXiv:2311.18193.

% No DeYO mechanism or baseline comparison is made in the maintained paper.
\KBbibitem{lee2024deyo}{Lee et al.}{2024} J. Lee, D. Jung, S. Lee, J. Park, J. Shin, U. Hwang, S. Yoon.
Entropy is not enough for test-time adaptation: from the perspective of disentangled
factors. \emph{ICLR}, 2024. arXiv:2403.07366.

% No comparison to this reset-policy analysis is made in the maintained paper.
\KBbibitem{lim2026reset}{Lim et al.}{2026} T. Lim, J. Hwang, K. Lee. When and where to reset matters for
long-term test-time adaptation. \emph{ICLR}, 2026. arXiv:2603.03796.

% No Tempora temporal-utility experiment or comparison is reported.
\KBbibitem{tempora2026}{Sreeram et al.}{2026} S. Sreeram, Y. D. Kwon, C. Mascolo.
Tempora: Characterising the time-contingent utility of online test-time adaptation.
\emph{ICML}, 2026. arXiv:2602.06136.

% No NTQR/error-correlation result is invoked in the maintained paper.
\KBbibitem{corradaemmanuel2024}{Corrada-Emmanuel}{2023} A. Corrada-Emmanuel. The logic of NTQR evaluations of
noisy AI agents: Complete postulates and logically consistent error correlations.
2023. arXiv:2312.05392.

% The two cited 2010 papers support the specific current impossibility comparison.
\KBbibitem{bendavid2012hardness}{Ben-David and Urner}{2012} S. Ben-David, R. Urner. On the hardness of domain
adaptation and the utility of unlabeled target samples. \emph{ALT}, 2012.

% No human-deferral or fairness analysis is made in the maintained paper.
\KBbibitem{madras2018learntodefer}{Madras et al.}{2018} D. Madras, T. Pitassi, R. Zemel. Predict responsibly:
improving fairness and accuracy by learning to defer. \emph{NeurIPS}, 2018.
